\section{Sistemi e impostazioni delle prove}
Tutti i sistemi ricevono gli stessi circuiti e operano nello stesso insieme
di cinque dispositivi: IBM Falcon 27, IBM Falcon 127, IBM Heron 133,
IBM Heron 156 e Quantinuum H2-56. Si usano le proprietà sintetiche congelate
di MQT Bench. Una scelta è ammessa solo se il dispositivo può ospitare
il circuito. Il processo di compilazione ha un limite di 100 secondi,
che comprende anche il suo avvio. Un fallimento conclude quel caso:
non si cerca una nuova configurazione per migliorarne l'esito.

Per le tre varianti LLM e per Random, la compilazione usa Qiskit con seme 0 e un solo
processo. La scelta è limitata a un catalogo di dodici configurazioni.
Le opzioni combinano livelli di ottimizzazione 2 e 3 con diverse modalità
di disposizione dei qubit e instradamento delle operazioni.
Il catalogo è lo stesso per questi quattro sistemi.

\subsection{LLM + RAG}
Questo sistema usa Qwen3.5-4B con pesi Q8\_0 e temperatura 0, nella
configurazione scelta sulla validation. Riceve le caratteristiche del circuito,
i dispositivi compatibili, le dodici configurazioni Qiskit e cinque esempi
recuperati dal solo train. L'assunzione è che circuiti con caratteristiche
simili possano fornire indicazioni utili sulla compilazione.

Il recupero confronta 49 caratteristiche numeriche con distanza Manhattan.
La trasformazione e la scala delle caratteristiche sono ricavate solo dal
train. L'indice usa 396 circuiti distinti per contenuto: i 26 duplicati
fra i 422 file train non diventano esempi aggiuntivi.
Il Dataset del RAG contiene compilazioni Qiskit già valutate; non contiene
i punteggi dei circuiti Test. Al modello viene fornita una rappresentazione
compatta delle informazioni, senza il codice QASM completo.

La risposta deve indicare una coppia dispositivo/configurazione e uno o due
fatti verificabili. Sono consentite al massimo tre risposte complete,
quindi due richieste di correzione. La prima scelta accettabile diventa
definitiva e viene compilata una sola volta. All'ultimo tentativo una coppia
ammessa può essere accettata anche se i fatti non sono tutti verificati:
questa condizione viene registrata. La spiegazione libera non è certificata.

Il limite di risposta è 4096 token e quello di attesa è 3600 secondi per
chiamata. Il contesto richiesto è 60000 token, la cache usa Q8\_0,
i gruppi di elaborazione sono 512 e 128 token e tutti gli strati del modello
sono richiesti sulla GPU. Il pensiero esteso è disattivato.
Il seme è 20260913; i parametri di generazione sono
\texttt{top\_p}=0,95, \texttt{top\_k}=40, \texttt{min\_p}=0,
penalità di ripetizione 1 e penalità di presenza e frequenza 0.
Non sono previste ripetizioni automatiche per problemi di trasporto.

\subsection{LLM no RAG}
Si usa lo stesso Qwen3.5-4B, con gli stessi pesi e tutti i parametri appena
descritti: temperatura 0, massimo 4096 token, contesto 60000 e lo stesso seme.
Restano disponibili caratteristiche del circuito, dispositivi compatibili
e dodici configurazioni, ma non vengono recuperati esempi.
Questo sistema permette di osservare quanto cambia la scelta quando
il modello dispone soltanto delle informazioni sul caso corrente.

La risposta deve riportare un solo fatto sulla capacità in qubit del
dispositivo scelto. Restano il massimo di tre tentativi, i controlli sulla
coppia e la stessa regola di accettazione finale dei fatti.
Il confronto con RAG comprende quindi sia gli esempi sia la diversa verifica
dei fatti richiesti; non isola esclusivamente l'effetto della loro presenza.

\subsection{MQT Predictor}
MQT Predictor 2.4.0 combina due componenti. Un classificatore supervisionato
sceglie il dispositivo; una politica di apprendimento per rinforzo (RL)
specifica del dispositivo decide i passaggi di compilazione.
Usiamo l'architettura ML seguita da RL descritta nel lavoro del 2025.
Il Training set per il selettore e il Dataset del RAG sono quindi risorse
diverse, costruite per scopi diversi.

Le cinque politiche RL sono state preparate prima del Test, con un obiettivo
di 100000 passi di addestramento per dispositivo (100352 al completamento
dei gruppi di aggiornamento). Ottimizzano la metrica di fedeltà attesa.
Il Test esegue una sola selezione e una sola compilazione RL per circuito,
con limite esterno di 100 secondi. Il seme interno di MQT non è controllato.
MQT può costruire una sequenza di passaggi e non è vincolato alle dodici
configurazioni Qiskit degli altri sistemi.

%%MQT_DETAILS%%

\subsection{Random}
Random estrae con probabilità uniforme una coppia fra tutte quelle compatibili
di dispositivo e configurazione Qiskit. Il seme deriva da 20260921 e
dall'impronta del circuito, così da conservare la scelta effettuata.
La coppia estratta viene compilata una sola volta, con seme Qiskit 0
e limite di 100 secondi. Un fallimento non provoca una nuova estrazione.

Random non usa esempi, modelli linguistici o addestramento.
Fornisce un riferimento semplice: permette di capire quanto la scelta
ragionata migliori rispetto a una scelta casuale nello stesso catalogo.
Una sola estrazione per circuito, tuttavia, non misura tutte le possibili
prestazioni del metodo casuale.


\subsection{LLM + Random RAG}
Questa variante mantiene modello, prompt, catalogo, parametri di generazione,
controlli e compilazione di LLM + RAG. Cambia il recupero: estrae cinque
esempi train uniformemente e senza reinserimento fra gli stessi candidati
compatibili. Non calcola la distanza Manhattan. Gli alias E1--E5 seguono
l'ordine di estrazione. Il seme del recupero è 20260927; quello per circuito
dipende anche dall'impronta del QASM. Il seme di generazione LLM resta 20260913.

La scelta della coppia dispositivo/configurazione è ancora affidata al modello.
Random, invece, estrae direttamente la coppia senza usare un LLM.
La campagna casuale conserva un contratto separato ed è stata eseguita
dopo aver osservato il Test originale. Non è una conferma indipendente:
un solo seme non misura la variabilità degli insiemi di esempi estratti.
Le differenze di tempo possono dipendere anche dal carico e dalla cache
nelle diverse sessioni.

I fatti dichiarati vengono sottoposti agli stessi controlli del RAG.
Una compilazione valida può comunque essere accettata con fatti non verificati.
Le ipotesi testuali non sono verificate e non vanno considerate affermazioni vere.
